\MakeAbstract{%
渗透测试通过模拟攻击者行为来发现系统漏洞，是Web安全防护中不可缺少的环节。随着大语言模型（LLM）推理能力的增强，已有多项工作尝试用LLM Agent来实现渗透测试的自动化。AutoPT是其中比较有代表性的一个，它用渗透测试状态机（PSM）对测试流程建模，约束了Agent的决策空间，但在实际使用中仍受限于状态转移粒度较粗、缺少浏览器交互工具以及只有命令行界面等问题。

本文针对这三个问题改进了AutoPT。本文设计了一套ReAct提示词模板由所有Agent共用，但各自有不同的角色边界。同时将AutoPT 的Scan阶段拆分为端口发现、服务知识匹配、漏洞扫描和结果评估四个子阶段，配合上下文压缩和Check结果匹配收窄等机制。为了加强对不同漏洞的自适应能力，本文新增了三个浏览器扩展工具用于表单交互，ReAct解析器也改造为多级容错机制。此外还搭建了一套基于Flask和JavaScript的Web Console管理界面。

实验使用了与原论文相同的20个CVE漏洞环境和三种模型，共进行300次独立测试。改进版总成功率为63.67\%，比原版高出约39\%；原版60组测试中有38组实验成功率为0\%，改进版压缩到了3组。弱模型GPT-3.5-turbo的变化最为突出，成功率从原版的约7\%跃升至54\%。同样以原论文中表现最好的GPT-4o-mini来计算成本，改进版单次成本降低69\%，平均耗时缩短42\%。这些数据说明，在不更换更强模型的前提下，对PSM架构内的流程控制和工具能力做针对性改进，就足以带来成功率和成本上的明显改善。
}{自动化渗透测试，大语言模型，智能体，渗透测试状态机，ReAct框架}

\MakeAbstractEng{
This thesis focuses on penetration testing, that is finding vulnerabilities in a system by simulating the behavior of an attacker. With the improving reasoning capability of large language models (LLMs), a number of programs have been developed to automate this process, and AutoPT is one of the most representative. AutoPT models the testing procedure with a penetration testing state machine (PSM) so that the agent does not drift off course. In practice, however, AutoPT is still limited by state transitions that are not fine-grained enough, the lack of browser interaction tools, and a command-line-only interface.

This thesis reworks AutoPT on all three fronts. A shared ReAct prompt template gives every agent the same output format while assigning each one a distinct role. Scan is no longer a single monolithic step; it now proceeds through port discovery, service knowledge matching, vulnerability scanning, and result evaluation in sequence. Context compression keeps token usage from growing excessively across ReAct iterations, and the Check phase only consults real tool output when deciding success or failure. Three new browser tools (fill\_element, select\_option, wait\_for\_selector) handle the form-based interactions that the original Playwright toolkit could not, and the ReAct parser has been rebuilt to tolerate the format deviations that chain-of-thought models tend to produce. In addition, a Web Console management interface based on Flask and JavaScript replaces the command-line interface.

We ran 300 independent tests (20 real CVE environments $\times$ 3 models $\times$ 5 runs each) under exactly the same setup as the original AutoPT paper. The improved version raises the overall success rate from roughly 25\% to 63.67\%. Among the 60 model--vulnerability combinations, 38 achieved a 0\% success rate in the original version, whereas only 3 still do in the improved one. The improvement is most pronounced for the weak model GPT-3.5-turbo, whose success rate rises from about 7\% to 54\%. Taking GPT-4o-mini, the best-performing model in the original paper, the improved version cuts the cost per test by 69\% and shortens the average runtime by 42\%. None of these gains required a more powerful model; they came from tightening process control, expanding the tool set, and improving usability.
}{automated penetration testing, large language model, agent, penetration state machine, ReAct framework}
